\chapter{Introduction}
\label{ch:introduction}

\section{Motivation and application context}
\label{sec:introduction:motivation}

Industrial facilities contain spatially distributed assets, restricted areas,
and conditions in which a potential fire must be assessed quickly. Aerial
robots can extend the field of view of a monitoring system and can approach
areas that are inconvenient or unsafe for a person to inspect. Their use does
not remove the sensing problem, however. An RGB camera records visual
appearance rather than temperature or combustion products, and an airborne
computer must process that stream within strict limits on memory, computation,
power, heat dissipation, payload and wireless connectivity.

Sending every camera frame to a remote server would make hazard detection
dependent on link capacity and remote availability. Edge inference places the
detector beside the camera and permits a compact event record to be transmitted
instead. Published work on edge-oriented UAV fire detection consequently
assesses predictive quality together with runtime and energy or resource cost
\cite{Vazquez2026EdgeFire}. Research on collaborative UAV edge intelligence
similarly treats computation, storage and link availability as joint design
constraints rather than assuming an unlimited central processor
\cite{Qu2024ElasticUAV}. These studies motivate local and cooperative
processing, but their numerical results cannot be transferred to the hardware,
data or two-stage decision process evaluated in this thesis.

This work studies a heterogeneous two-vehicle system. The \emph{Sentinel}
performs initial visual hazard detection on an NVIDIA Jetson Nano. The
\emph{Scout} performs secondary visual hazard verification on a ModalAI VOXL~2.
The design aims to reserve the Scout for candidates that persist at the
Sentinel and to preserve a traceable relationship among the initial candidate,
the verification request, the returned verdict and the final decision. The
title's phrase ``industrial fire detection'' identifies the target application
domain; it does not imply validation at an operating industrial facility. The
final physical target is visible flame, while broader industrial sensing and
physical smoke verification remain outside the evaluated scope.

\section{Problem statement}
\label{sec:introduction:problem}

A single detector on a constrained UAV faces competing objectives. Lower model
capacity or resolution may reduce runtime cost but can remove visual detail.
A permissive threshold may retain weak hazards but can also pass fire-like
backgrounds. Requiring persistence across frames can suppress transients, but
it can delay or discard short events. A second verifier can reject some
candidates, but it introduces another model, camera, computer, communication
path and timeout condition. The architectural benefit is therefore an
experimental question rather than an automatic consequence of adding a second
vehicle.

The problem addressed here is to determine how far a resource-aware division
of responsibilities can support visual fire detection and verification on two
heterogeneous edge platforms. The study must compare detector choices, measure
the selected runtimes on the devices that execute them, examine sustained
operation with concurrent services, and determine what the temporal and tiered
verification stages do to event retention and decision time. It must also test
the two-vehicle mission logic under defined simulation conditions and examine
whether predictive performance transfers beyond the development data.

This formulation deliberately separates an engineering prototype from an
operational safety system. RGB predictions do not provide physical fire ground
truth. Secondary verification by another RGB model is an independent visual
assessment, not proof of combustion. The laboratory system is stationary, and
the simulation uses modeled vehicles and scenario coordinates. The thesis
therefore evaluates a bounded research prototype and records the evidence
required before broader deployment claims could be considered.

\section{Terminology and system interpretation}
\label{sec:introduction:terms}

Three terms in the title require precise interpretation.

\begin{description}
  \item[Tiered.] Detection and verification responsibilities are assigned to
  distinct, ordered stages across the two edge platforms. The Sentinel detects
  a candidate and decides whether to request verification; the Scout returns an
  independent verdict; the Sentinel records one correlated terminal decision.
  The term describes role specialization and decision flow; it does not imply
  a multi-level command organization or feature-level neural-network fusion.

  \item[Cooperative UAV pair.] The prototype contains two cooperating aerial
  platforms that exchange typed messages over a shared network. This is the
  smallest multi-vehicle configuration needed to study role specialization. It
  does not support conclusions about scaling to a larger fleet or swarm.

  \item[Industrial.] Industrial describes the intended monitoring context and
  its resource and safety constraints. The retained evidence comes from
  offline datasets, exact-device measurements, stationary laboratory trials
  and simulation. No industrial-site deployment was conducted.
\end{description}

Throughout the thesis, \emph{hazard detection} means the Sentinel's initial
visual identification and \emph{hazard verification} means the Scout's
secondary assessment. The \emph{temporal filter} is the consecutive-frame rule
applied before an alert is published. The title's \emph{tiered edge AI
architecture} denotes this division of responsibilities between the two edge
platforms and the ordered detection--verification decision flow. ``Fire'' in
the final live evaluation refers to a visible flame.
Historical fire/smoke artifacts are reported only as development history or
offline evidence; physical smoke verification is outside the validated live
scope.

\section{Research questions and objectives}
\label{sec:introduction:rqs}

The investigation is organized around five research questions.

\begin{description}
  \item[RQ1.] What predictive-performance and resource trade-offs do the
  candidate detectors exhibit, and which are supported for each edge platform?

  \item[RQ2.] Can the selected configurations operate sustainably with the
  required concurrent services?

  \item[RQ3.] How do temporal filtering and Scout verification affect event
  retention, false alerts and decision latency?

  \item[RQ4.] Does the implemented two-vehicle mission behave correctly and
  handle faults in the tested simulation?

  \item[RQ5.] How does predictive performance transfer to eligible external
  UAV and industrial-like datasets?
\end{description}

Table~\ref{tab:rq-objectives} converts these questions into objectives and
identifies where each is answered. An unanswered question remains an evidence
gap; it is not assigned a favourable result by inference.

\begin{table}[H]
  \centering
  \caption{Research-question objectives and thesis traceability.}
  \label{tab:rq-objectives}
  \small
  \begin{tabularx}{\textwidth}{@{}p{0.10\textwidth}X p{0.22\textwidth}@{}}
    \toprule
    Question & Operational objective & Principal chapters \\
    \midrule
    RQ1 & Compare predictive metrics, artifact fidelity and exact-device
    resource behaviour without mixing laptop and edge measurements. &
    Chapters~4, 6--8 \\
    RQ2 & Evaluate the selected configurations with the camera, communication
    and logging services required by the workflow over a defined duration. &
    Chapters~5--8 \\
    RQ3 & Compare single-frame, temporally filtered and verified decisions on
    admissible labelled units, including the latency introduced by each stage. &
    Chapters~5--8 \\
    RQ4 & Exercise nominal mission branches and preregistered faults in the
    simulator after its readiness gates are satisfied. & Chapters~5--8 \\
    RQ5 & Evaluate frozen models on external data admitted through provenance,
    overlap and annotation checks. & Chapters~4 and 6--8 \\
    \bottomrule
  \end{tabularx}
\end{table}

\section{Contributions}
\label{sec:introduction:contributions}

The thesis makes four evidence-bounded contributions. First, it specifies a
two-node tiered edge architecture in which a recall-oriented Sentinel
stage and a precision-oriented Scout stage exchange compact, correlated
decision records. Second, it develops an auditable protocol with explicit
experiment, trial and event identities; monotonic timing; model hashes;
timeouts; terminal decisions; and append-only logging. Third, it establishes a
comparison framework that connects dataset-controlled detector screening to
runtime feasibility on the intended devices while preserving the distinction
between historical and final deployment artifacts. Fourth, it integrates the
decision process with a safety-constrained two-vehicle simulation architecture
and defines gates that prevent development demonstrations from being reported
as formal reliability evidence.

These contributions concern architecture, implementation and evaluation
method. Claims about effectiveness are made only where retained measurements
support them. In particular, the study does not claim universal reliability,
zero false positives, autonomous industrial flight, geographic hazard
localization or physical smoke verification.

\section{Scope and evidence boundaries}
\label{sec:introduction:scope}

The main evidence categories are kept separate. Offline predictive evaluation
tests frozen image splits. Exact-device experiments measure particular model
artifacts and runtime configurations. Stationary laboratory trials demonstrate
the integrated physical detection--verification path under recorded
conditions. Simulation evaluates mission software under modeled dynamics and
scenario inputs. External-data evaluation, when its admission conditions are
met, concerns transfer to those datasets only. A result in one category is not
presented as evidence in another.

The physical evaluation uses stationary vehicles with their propellers
removed. The final live hazard is a handheld visible flame. Image-space box
coordinates describe the Sentinel camera observation; they are neither
geographic coordinates nor a position in the Scout camera frame. Integrated
AI-controlled flight, common-world localization, physical-smoke trials,
industrial-site trials and systems containing more than two vehicles are
outside the validated scope.

The deployed Sentinel ROS workspace contains a 640-pixel TensorRT engine. An
earlier benchmark used a different 416-pixel engine. Their input and output
shapes and artifact identities differ, so measurements from the earlier engine
are not treated as performance measurements of the final deployment. The same
identity rule applies to every model and runtime comparison in later chapters.

\section{Thesis organization}
\label{sec:introduction:organization}

Chapter~2 reviews visual hazard detection, lightweight detectors, edge
optimization, cooperative perception, temporal processing and ROS~2/DDS.
Chapter~3 derives the requirements and presents the tiered architecture.
Chapter~4 documents dataset preparation, detector development and deployment
artifacts. Chapter~5 describes the software, device and communication
implementation. Chapter~6 defines the experimental and statistical methods,
and Chapter~7 reports accepted results. Chapter~8 answers the research
questions and examines limitations and threats to validity. Chapter~9 states
the evidence-supported conclusions and the work required for broader
validation.
